%%%% TwoCascade — SURS poster, Better Poster format
%%%%
%%%% Base template: Rafael Bailo, betterposter-latex-template (GNU GPL v3.0)
%%%%   https://github.com/rafaelbailo/betterposter-latex-template
%%%% Original design: Mike Morrison, https://twitter.com/mikemorrison
%%%%
%%%% TwoCascade adaptation — Gary Mei, Georgia Tech ISyE (SURS).
%%%% Directive: okf/decisions/d-038-advisor-poster-two-comparisons.md
%%%% Slot plan and evidence provenance: okf/poster/poster-plan.md
\documentclass[fleqn]{betterposter}

\usepackage{geometry}
\usepackage{tikz}
\usepackage{enumitem}
\geometry{paperwidth=36in, paperheight=48in}

\usepackage{xcolor}

%%%%%%%%%%%%%%%%%%%%%%%%%%%%%%%%%%%%%%%%%%%%%%%%%%%%%%%%%%%%%%%%%%%%%%%%%%%%%%%
%% PALETTE — matches the advisor-update site and the .lavish artifacts
%%%%%%%%%%%%%%%%%%%%%%%%%%%%%%%%%%%%%%%%%%%%%%%%%%%%%%%%%%%%%%%%%%%%%%%%%%%%%%%
\definecolor{tcCream}{HTML}{F6F3EA}   % page / side-column background
\definecolor{tcInk}{HTML}{1E2530}     % body text
\definecolor{tcInkSoft}{HTML}{5B6672} % secondary text
\definecolor{tcTeal}{HTML}{2E6E76}    % cool accent — structure, geometry
\definecolor{tcRust}{HTML}{BD5A2E}    % warm accent — fear
\definecolor{tcRule}{HTML}{D8D2C2}    % hairlines

%% Third CATEGORICAL colour, for the GIRG family in the hero figure.
%% Not tcInkSoft: grey reads as "disabled/secondary" next to two saturated hues,
%% and grey-vs-teal was the pair most likely to collapse at viewing distance.
%% This is the Okabe-Ito reddish-purple, darkened for contrast on cream --
%% chosen because purple is the standard third that stays separable from both
%% orange and blue-green under deuteranopia and protanopia.
\definecolor{tcPlum}{HTML}{8E4A72}    % third family — geometry (GIRG)

%% Side columns: cream ground, ink type.
\renewcommand{\columnbackgroundcolor}{tcCream}
\renewcommand{\columnfontcolor}{tcInk}

%% Main column: deep teal ground, cream type.
\renewcommand{\maincolumnbackgroundcolor}{tcTeal}
\renewcommand{\maincolumnfontcolor}{tcCream}

\begin{document}
\betterposter{
%%%%%%%%%%%%%%%%%%%%%%%%%%%%%%%%%%%%%%%%%%%%%%%%%%%%%%%%%%%%%%%%%%%%%%%%%%%%%%%
%%%%%%%% MAIN COLUMN — headline and core finding visual
%%%%%%%%%%%%%%%%%%%%%%%%%%%%%%%%%%%%%%%%%%%%%%%%%%%%%%%%%%%%%%%%%%%%%%%%%%%%%%%

\maincolumn{
%%%% Main space: Headline INFORMS, figure below SHOWS the finding.
%%
%% Headline rewritten by Gary 2026-08-06 (was the "Who you are connected
%% to..." construction): states the contribution (fear as a new failure
%% route) and what the poster shows (its interaction with heterogeneity
%% and geometry).
%%
%% Baselineskip was 0, which is why the two rendered lines had that large uneven
%% gap -- TeX falls back to \lineskip when \baselineskip cannot be honoured.
{\fontsize{80}{80}\selectfont
\textbf{Fear} is a new way for nodes to fail.
\textbf{This} is how fear interacts with heterogeneity and geometry.
\par}% <-- CLOSED FONT BRACE HERE BEFORE THE FIGURE!
\vspace{50pt}

%% [REAL FIGURE — Comparison 1, heterogeneity (ER vs CM), n=10000, 500 trials/pt]
%%
%% Source: results/figures/poster_comparison.png, regenerable from
%% results/processed/poster_comparison.json (hero data 301d560; departure
%% analysis now 8-point after the mu-bar 0.5/0.6 densification arms,
%% 0e739a1/d25f9f2; pipeline blind-gated AUDIT PASS S-066). Legend and fear
%% encoding are baked into the PNG.
%%
%% GIRG joined this axis 2026-08-02 (C3b production run, 500 trials/cell,
%% calibrated w_min=0.186377): Gary accepted the 9-curve figure as the hero,
%% resolving the previously-open one-axis-vs-separate question. The two D-038
%% comparisons stay separable BY READING: CM-vs-ER = heterogeneity,
%% GIRG-vs-CM = geometry. Record lives in the C3b commit chain.
%%
%% Full-bleed hero (Gary, 2026-08-04): width=1\linewidth with a 1\textheight
%% ceiling, keepaspectratio -- the width constraint binds; the legend +
%% Sec 5.4 engine-disclosure row below is pulled up into the freed
%% whitespace with the negative vspace.
%% DPI: all poster figures verified 300 dpi (S-066; the old "regenerate
%% before print" warning here was stale).
\includegraphics[width=1\linewidth, height=1\textheight, keepaspectratio]{img/poster_comparison}\\

%% ⚠ §5.4 DISCLOSURE — RESOLVED by D-047 (Gary, 2026-08-06): the printed
%% engine-split sentence is DELIBERATELY OMITTED from the poster (audience
%% legibility — "other people in the audience won't get it"). This
%% overrules the S-066 audit's carried condition with a recorded decision;
%% the disclosure survives HERE and in okf/poster/presentation-script.md
%% (spoken if asked). The hero mixes engines: ER on C++, power-law on
%% Python, GIRG on both (cross-engine agreement verified, §5.4).
\vspace{-250pt}
\begin{minipage}[c]{0.55\linewidth}
    \centering
    \includegraphics[width=1\linewidth, height=0.08\textheight, keepaspectratio]{img/poster_comparison_legend}
\end{minipage}\hfill
\begin{minipage}[c]{0.42\linewidth}
    \raggedright
    \fontsize{35}{30}\selectfont % <-- MANUAL SIZE DIAL


    Measured at matched mean degree $\langle k\rangle = 4.533$, $n = 10000$.

    500 trials per point.

\end{minipage}

}{
%%%% Bottom space of center column left intentionally clean.

}
}{
%%%%%%%%%%%%%%%%%%%%%%%%%%%%%%%%%%%%%%%%%%%%%%%%%%%%%%%%%%%%%%%%%%%%%%%%%%%%%%%
%%%%%%%% LEFT COLUMN — the argument, in reading order (Gary's 2026-08-06
%%%%%%%% structure: existing graph families -> contribution -> fear
%%%%%%%% definition -> what we look for; QR lives in the RIGHT column; the
%%%%%%%% D-046 model schematic lives on the QR demo site's landing screen,
%%%%%%%% not on the poster).
%%%%%%%%%%%%%%%%%%%%%%%%%%%%%%%%%%%%%%%%%%%%%%%%%%%%%%%%%%%%%%%%%%%%%%%%%%%%%%%
\title{Confidence \\
Augmented\\
Bootstrap \\
Percolation\vspace{-0.4em}}  % <-- Pulls your name up closer to the title!

\author{Gary Mei\vspace{.5em}}

\institution{Georgia Tech ISyE - SURS\\
Advisor: Prof.\ Souvik Dhara}

%% Place \raggedright HERE, after the header block, to wrap body text cleanly!
\raggedright

%% Pulls the first section up closer to the advisor line!
\vspace{-2.0em}

\section{Existing Work}
\vspace{-0.3em}
\begin{enumerate}
    %% \H{o}/-- macro forms, not UTF-8 literals: the bold sans font
    %% (ec-lmssbx10) has no precomposed ő or en-dash glyph and silently
    %% drops them (rendered "ErdsRényi" -- caught in the S-066 verify pass).
    \item \textbf{Erd\H{o}s--R\'enyi graph:} Independent edge probability $p$ between any two nodes.
    \item \textbf{Power law configuration graph:} A few nodes (aka hubs) hold most edges; most nodes have very few.
    \item \textbf{GIRG (Geometric Inhomogeneous Random Graph):} Power-law plus geometry; closer nodes are likelier to connect.
\end{enumerate}

\vspace{-1.5em}

\section{My contribution}

\vspace{-0.3em}
A node now activates if \textbf{either}
condition is met:
\vspace{0.2em}
% Converted equations to inline math ($...$) to save ~2 inches of vertical space!
\begin{enumerate}[leftmargin=*, nosep, itemsep=0.3em]

\item \text{Percolation:}
At least $r$ of its neighbors have failed ($|\{j \sim i : j \in A_t\}| \geq r$).

\item \textbf{Fear:}
A Bernoulli variable with probability $f_ig_t$ that a node activates ($U_i < f_i\, g_t$).

\end{enumerate}

\vspace{-1.2em}
\section{How is fear defined?}
\vspace{-0.3em}
%% g_t wording (S-066 fix): the engine's fear field is WINDOWED -- the
%% average failed share over the last window_len=5 generations (uniform
%% weights), NOT last generation alone. Keep the printed definition honest.
\begin{enumerate}
    \item $f_i \in [0,1]$: Node $i$'s fear susceptibility (1 = highest).
    \item $g_t$: the share of nodes that failed recently --- averaged over
    the last 5 generations.
\end{enumerate}


\vspace{-1.5em}
\section{What We Look For}
\vspace{-0.3em}
Holding $n$ and $\langle k\rangle$ fixed, we isolate three effects:
\vspace{0.2em}
\begin{itemize}[leftmargin=*, nosep, itemsep=0.3em]

  \item \textbf{Heterogeneity:}

  homogeneous degrees vs.\ heavy-tailed hubs.

  \item \textbf{Geometry:}

  pure random links vs.\ local spatial clustering.

  \item \textbf{Fear:}

  baseline percolation vs.\ confidence augmentation.

\end{itemize}

\vfill

%% [QR box lives at the bottom of the RIGHT column (Gary's 2026-08-06
%% layout).]

}{
%%%%%%%%%%%%%%%%%%%%%%%%%%%%%%%%%%%%%%%%%%%%%%%%%%%%%%%%%%%%%%%%%%%%%%%%%%%%%%%
%%%%%%%% RIGHT COLUMN — supporting evidence and the ISyE recommendation.
%%%%%%%%%%%%%%%%%%%%%%%%%%%%%%%%%%%%%%%%%%%%%%%%%%%%%%%%%%%%%%%%%%%%%%%%%%%%%%%
\raggedright

\vspace{-1.9em}

\section{Fear is an amplifier with a dial.}
%% PANEL SWAP (Gary, 2026-08-06): this slot now shows
%% results/figures/famcompare_probability.png -- P(systemic | mu-bar) at the
%% bounded seed a=2 for power-law/GIRG/ER, mu-bar 0..1.0
%% (scripts/plot_famcompare_probability.py, ext2 arms mu-bar {0.8,0.9,1.0}
%% run 2026-08-06, base_seeds 20260806/07/08, 500 trials/cell, provenance
%% guards extended). The famcompare_ratio multiplier figure and its >=26x /
%% 6.6x / 8.3x claims moved OFF the poster.
%% ⚠ EVIDENTIARY TIER: the probability panel + ext2 cells are post-S-066-gate,
%% script-level self-checked (guard-protected, not blind-reviewed). The
%% S-066 AUDIT PASS covers the hero + Fear-and-Structure pipelines and the
%% (now unprinted) ratio pipeline.
%% PRESENTER NOTES (kept for the talking script; the multiplier framing is
%% still true, just no longer printed):
%%   Fear's multiplier on the heavy-tailed families climbs with the dial:
%%   ~1.3x (power-law; GIRG is 0.8x there, below the no-effect line with a
%%   wide CI [0.35, 1.75] -- direction-only) at mu-bar=0.1, ~3x at 0.4,
%%   ~6.6x (power-law) / 8.3x (GIRG) at 0.7 -- none censored.
%%   Amplification depends on structure. With a bounded seed (a=r=2), ER
%%   stays at 0/500 through mu-bar=0.7, first ignites at 0.8 (9/500), then
%%   7.6% at 0.9 and 25.6% at 1.0: fear amplifies cascades the structure
%%   makes possible; only near-certain panic lets it start one on ER.
%%   GEOMETRY NOTE: on the matched ensemble GIRG and power-law INTERLEAVE
%%   with overlapping CIs throughout -- geometry doesn't measurably change
%%   fear's effect in either direction.
%\textbf{Fear is an amplifier with a dial.}
%\vspace{0.2em}

When we set $a = 2$, the Erd\H{o}s--R\'enyi graph doesn't activate at all
until $\bar\mu = 0.8$ --- and even then in only $\sim$2\% of trials. The
effects of fear on the GIRG and power-law model appear relatively equal.

\begin{center}
\makebox[\linewidth][c]{% Forces centering even when wider than margins
  \fcolorbox{tcRule}{tcCream}{%
    \begin{minipage}[t]{1.2\linewidth} % Increased to 120% width
      \centering
      %% S-066: 0.70, square figure, width=height; sized so the column tail
      %% (bullet 3 + Next Steps + QR) fits on-page.
      \includegraphics[width=1\linewidth, height=20in, keepaspectratio]{img/famcompare_probability}
    \end{minipage}%
  }%
}
\end{center}

\vspace{-1.9em}

\section{Fear and Structure}
%% 2026-08-03: replaces the D-012 measured/predicted-RATIO framing (see
%% okf/decisions, D-012 poster evidentiary choices). (1-mu_bar)^2 is derived
%% for ER only; scoring CM/GIRG's measured crossing against it treated an
%% ER-only formula as a baseline they were never predicted to meet. This
%% section instead states each family's own measured percentage decrease
%% against ITS OWN mu_bar=0 anchor -- self-contained per family, and the
%% ER-only theory now appears only as a labeled, dashed reference. mu_bar=0.7
%% re-enters for every family (its old exclusion was an artifact of the
%% predicted-crossing floor check, not a property of the measurement itself);
%% the figure instead flags points within 2.2x of the true structural floor
%% a=r=2, where the floor -- not fear -- may be doing some of the compressing.
%% Numbers below: verify from results/processed/poster_fear_structure.json,
%% don't hand-copy. Pipeline blind-gated AUDIT PASS S-066.
%% 2026-08-04 (Gary): the ER-only (1-mu_bar)^2 dashed reference is REMOVED
%% from the figure and the y-axis is inverted so a decrease reads downward.
%% Condensed 2026-08-04 (Gary): all analysis above the figure, nothing below.

The same fear that cuts ER's tipping point by 65\% cuts the heavy-tailed
networks' by only $\sim$35\% at $\bar\mu = 0.4$. Hubs blunt fear's
effect, and geometry doesn't change that.

%% nu sentence kept off the print (Gary; ER-only claim, lives in the
%% presentation script):
%Fear also sharpens the transition faster as the network grows: $\nu(\mu = 0) = 5.61 \pm 0.18 \to\nu(\mu = 0.3) = 4.82 \pm 0.15$, verified.

\vspace{1em}

\begin{center}
\makebox[\linewidth][c]{% Forces centering even when wider than the margins
  \fcolorbox{tcRule}{tcCream}{%
    \begin{minipage}[t][6.5in][c]{1.2\linewidth} % Increased to 120% width
      \centering
      \includegraphics[width=1\linewidth, keepaspectratio]{img/poster_fear_structure}
    \end{minipage}%
  }
}
\end{center}

\vspace{0em}

% \section{How To Intervene}
%% INTERVENE PANEL CUT (Gary, 2026-08-06 final revision) -- section + the
%% three printed AUDIT-PASS anchors + the "Next:" line removed from print
%% to make room for the large QR box below. All of it lives on as SPOKEN
%% material in okf/poster/presentation-script.md.
%% TALKING POINTS:
%%   In a network, fear near an active failure is harmless, because those
%%   nearby nodes would likely fail anyway. The real danger is fear reaching
%%   distant nodes and starting new chain reactions --- so intervention
%%   should prioritize regions far from the crisis, not just the active
%%   crisis zone.
%%   Printed-anchor claims (all AUDIT-PASS: ignition gate C-Q4(iii); tilt
%%   monotonicity C-Q4(i)/Task X; local-vs-global dichotomy C-Q5(iii)):
%%   cap the concentration (tau>3: 0 systemic in 16,000 trials); calm the
%%   hubs (beats periphery at every fear level); contain the news (local
%%   fear ~ no fear). NOTE tau = degree-distribution tail exponent, not r.

%% QR box (bottom-right, Gary's 2026-08-06 layout, stacked + large). The
%% demo's landing screen opens on the model schematic.

\vspace{.5em}

\begin{center}
\fcolorbox{tcRule}{tcCream}{%
  \begin{minipage}{1\linewidth}
    \vspace{0.5em}
    \centering % Centers both the QR code and the text

    % 1. THE QR CODE
    % Changed from taking up 23% of the box to 50% (Dial this up to 0.7 or 0.8 if you want it huge)
    \includegraphics[width=0.7\linewidth]{img/qrcode}

    \vspace{0.5em} % Adds a little breathing room between the image and text

    % 2. THE TEXT
    \color{tcInk}%\bfseries
    \fontsize{30}{38}\selectfont % Dialed up from 30/36 to 40/48!
    Scan to interact
    with live graphs

    \vspace{0.5em}
  \end{minipage}%
}
\end{center}

}
\end{document}
